\begin{claim}{C.7g}{Proved}
Let $K/\Q_p$ be finite, put $F=K\Q_p(\mu_\infty)$, and let $F'/F$
be finite. On every smooth geometrically connected affine genus-one
curve over $F'$, ordinary point sections are injective. More generally,
an invariant path between two ordinary or tangential point sections
forces their underlying points on the smooth compactification to be
equal. In particular an ordinary point section cannot equal a
tangential section, and tangents at distinct cusps cannot equal.
No assertion distinguishing all directions at one cusp is intended.

Consequently, for any hyperbolic $U/K$, every element of the full
ordinary fixed group $\pi_1(U_{\bar K})^{s_x(G_F)}$ acts trivially
in the monodromy of every finite geometric cover whose smooth
compactification has genus zero or one. For a geometric Galois
cover of either genus, the image of a tangential full fixed group
is exactly the corresponding cusp inertia subgroup. These conclusions
hold for every open subgroup of $G_F$ and for every residue prime.
They include the algebraic tripod point $2$ and direction $1$.
They still do not determine the full group in the remaining kernel.
\end{claim}
\begin{proof}
We first prove a finite separation assertion on a complex elliptic
curve $E$. Let $\lambda:E(\C)\setminus\{0\}\to\R$ be its normalized
Green function, with respect to Haar probability measure $\mu$.
The only analytic properties needed are
\[
 \Delta\lambda=\mu-\delta_0,\qquad
 \lambda(z)=-\log|z|+\rho(z)\quad\hbox{near }0,               \tag{C.7g.1}
\]
where $\rho$ is smooth, and smoothness away from $0$.
For an exact primary description, Baker--Petsche, Section~2.2,
equation~(10) and Proposition~1, gives these properties
\cite[author pp.~5--6]{BC_BakerPetsche2005}.
Indeed their convergent product expression shows that subtracting
$-\log|z|$ leaves a smooth function near zero; differentiating it
gives
\[
              \nabla\lambda(z)=-z/|z|^2+O(1).              \tag{C.7g.2}
\]
Here the complex coordinate is identified with a real two-dimensional
coordinate. The choice of an overall positive Laplacian normalization
does not affect any separation argument below.

If $T,S$ are distinct torsion points, choose $m>0$ with
$m(T-S)=0$ and a nonzero rational function $f_{T,S}$ satisfying
\[
                         \operatorname{div}(f_{T,S})=m(T-S).
\]
Such a function exists because $E\simeq\operatorname{Pic}^0(E)$.
Put
\[
             \psi_{T,S}(Z)=\lambda(Z-T)-\lambda(Z-S).
\]
The function $\log|f_{T,S}|+m\psi_{T,S}$ has zero Laplacian.
Its logarithmic singularities cancel at $T$ and $S$, so it extends
to a harmonic function on the compact torus and is constant.
Thus, for $P,Q$ away from $T,S$,
\[
 \log\left|\frac{f_{T,S}(P)}{f_{T,S}(Q)}\right|
       =-m\bigl(\psi_{T,S}(P)-\psi_{T,S}(Q)\bigr).          \tag{C.7g.3}
\]

We claim that finitely many such functions $f_1,\ldots,f_r$ can
be chosen so that for every pair $P\ne Q$ in $E(\C)$, at least
one $f_i$ is regular and nonzero at both points and
$|f_i(P)|\ne|f_i(Q)|$. The regularity condition is part of the
claim, including when $P$ or $Q$ occurs in the support of another
chosen function.

First, the entire family separates every distinct pair. Otherwise,
for fixed $P\ne Q$, the continuous function
\[
             T\longmapsto\lambda(P-T)-\lambda(Q-T)
                                                               \tag{C.7g.4}
\]
would take a constant value on all torsion points outside $P,Q$.
Torsion is dense in the complex torus. Therefore the function would
be constant on $E(\C)\setminus\{P,Q\}$, which contradicts its
opposite logarithmic singularities at $P$ and $Q$. This proves
separation by some $\psi_{T,S}$ with support disjoint from $P,Q$.

Next fix any $Z\in E(\C)$ and a torsion point $S\ne Z$.
Choose two torsion centres $T_1,T_2$ sufficiently close to $Z$,
in two small cones whose directions are linearly independent over
$\R$. Equation (C.7g.2) and density of torsion show that
$d\psi_{T_1,S}(Z)$ and $d\psi_{T_2,S}(Z)$ are linearly independent:
the singular radial terms can be made arbitrarily large, whereas
the contribution from $S$ and the smooth remainder stay bounded.
The real inverse function theorem supplies a neighbourhood $B_Z$
of $Z$, disjoint from these three centres, on which the two
functions together are injective. Choose finitely many such
neighbourhoods covering $E(\C)$. The union of their products
$B_Z\times B_Z$ is an open neighbourhood of the diagonal in
$E(\C)^2$; on it the finitely many selected functions already
separate every unequal pair and are regular at the pair in question.

The complement of this neighbourhood is compact and disjoint from
the diagonal. The first separation argument gives, near each of
its points, one function regular at both coordinates and having
unequal absolute values. Finitely many of these neighbourhoods
cover the compact complement. Adding their functions proves the
finite separation claim. Let $S_0\subset E(\C)_{\rm tors}$ be
the finite union of all their zeros and poles, and choose a positive
integer $N$ annihilating $S_0$.

We apply this construction over a cyclotomic local field. Let
$V/F'$ be the given affine genus-one curve. After a further finite
extension of $F'$, choose one geometric point of the nonempty
boundary as the origin $O$ of its compactification $E$. This
makes $E$ an elliptic curve. Such an extension is harmless for
proving uniqueness: an invariant path remains invariant after
restriction. Every finite extension of $F$ is $M\Q_p(\mu_\infty)$
for some finite local $M/K$, and every finite collection of its
elements lies in a finite local extension. Thus all subsequent
finite algebraic data can be placed in such an $M$.

For clarity about complex embeddings, choose a finitely generated
field $k_0\subset\bar K$ over which $(E,O,V)$ is defined, and
let $L$ be its algebraic closure in $\bar K$. Choose an embedding
$L\hookrightarrow\C$. All complex torsion points of this
elliptic curve lie in $E(L)$, since each torsion group scheme is
finite and $L$ is algebraically closed. The functions in the
finite separation argument can therefore be taken in $L(E)$.
The selected finite support $S_0$ and integer $N$ are fixed before
considering the endpoints of a proposed invariant path.

Consider the isogeny covering
\[
 V_N=E\setminus[N]^{-1}(E\setminus V)
                  \ \xrightarrow{[N]}\ V.                 \tag{C.7g.5}
\]
It is finite etale and connected, and its omitted points include
$E[N]=[N]^{-1}(O)$, hence include $S_0$. Every $f_i$ is therefore
an invertible regular function on $V_N$.

Suppose that two ordinary or tangential sections on $V$ admit an
invariant path and have distinct underlying compactified points.
After placing the covering, all $f_i$, and its two lifted endpoints
in a finite local extension $M$, lift the path starting at any
chosen lift of the first geometric or tangential basepoint. Its
terminal lift has some underlying point $Q$, while the initial
one has underlying point $P$. They are distinct because their
images under $[N]$ are distinct. The lifted path is invariant
under $H_M=G_{M\Q_p(\mu_\infty)}$; this follows by uniqueness of
path lifting once the selected initial basepoint is fixed.
Tangential lifts exist without ramification in (C.7g.5), since
$[N]$ is etale on the entire compact elliptic curve.

Only a finitely generated extension of $L$, obtained by adjoining
the coordinates of $P,Q$, needs to be embedded in $\C$.
The chosen embedding of $L$ extends to it: choose a transcendence
basis in $\C$ over the countable field $L$, then extend across the
remaining algebraic extension. No embedding of the entire local
field preserving its valuation is used. The finite separation
claim applies to these complex images of $P,Q$, so choose an $f_i$
regular and nonzero at both endpoints with
\[
                              |f_i(P)|\ne|f_i(Q)|.           \tag{C.7g.6}
\]
If an endpoint is tangential, fill that boundary point before
applying $f_i$. This is legitimate because $f_i$ is regular and
invertible there: its induced $\mathbf G_m$-section is the ordinary
point section with value $f_i(P)$ or $f_i(Q)$, independently of the
nonzero tangent coefficient. The invariant path therefore implies
\[
        \kappa_p\bigl(f_i(P)/f_i(Q)\bigr)|_{H_M}=0.
\]
The multiplicative local restriction lemma of C.7b makes this
ratio a root of unity. Under every complex embedding its absolute
value is one, contradicting (C.7g.6).

This proves uniqueness of the underlying compactified point.
Ordinary distinct points, an ordinary point and a cusp, and two
distinct cusps have distinct underlying points, giving all three
assertions. We have not tried to distinguish tangent coefficients
at the same cusp, where the statement of the theorem deliberately
makes no claim.

Finally take a finite geometric cover of any hyperbolic $U/K$
whose compactification has genus one. It is affine, since a finite
etale cover of a proper hyperbolic curve has compact genus at least
two. The preceding result applies after any finite local extension
needed to define this cover and its endpoints. An invariant ordinary
loop on $U$ lifts from each point of the fibre to another point,
and uniqueness forces those two points to coincide. Thus the
entire monodromy permutation is trivial, whether or not the cover
is Galois. In the Galois tangential case, an image outside cusp
inertia would change the underlying lifted cusp, which has just
been excluded. Conversely the full cusp inertia is fixed on the
cyclotomic kernel and maps onto that inertia subgroup. The genus-zero
case is C.7d. The same argument after restriction handles every
open subgroup of $G_F$.
\end{proof}

The elliptic abelian restriction alone would not prove this statement.
For example, on
\[
       E:Y^2=X^3+1,\qquad t=(Y+1)/2:E\longrightarrow\mathbf P^1,
\]
the degree-three map is branched precisely at $0,1,\infty$.
Above $t=2$ lies $P=(2,3)$, with
$2P=(0,1)$ and $3P=(-1,0)$, so $P$ has order six.
If $\tau(X,Y)=(\zeta_3X,Y)$, then $P-\tau P$ is nonzero
two-torsion: explicitly $P-\tau P=(-\zeta_3^2,0)$, and doubling
it gives zero because $\tau$ fixes $2P$.
Thus distinct points of this fibre already have torsion Abel
difference. C.7g distinguishes their sections by the additional
units on the isogeny cover (C.7g.5); it does not replace a torsion
condition by equality without that step.

The displayed arithmetic is checked with exact rational arithmetic in
$\Q[\zeta_3]/(\zeta_3^2+\zeta_3+1)$ by
\path{code/BC_elliptic_detection_checks.py}; its eight checks concern
this example, while the finite separation theorem is proved above.
